
\section{Sheaves: }

\subsection{Presheaves, Sheaves, and Stalks:}


\begin{exercise}{2.2.E.} 
    Let $\FF(U)$ be the set of maps $U \to S$ that are locally constant, i.e. for any point $p \in U$, there is an open neighborhood of $p$ where the function is constant. Show that this is a sheaf. We denote this, the constant sheaf, by $\underline{S}$.
\end{exercise}
\begin{solution}
    This is clearly a presheaf with the usual restriction on functions. Thus, we only need to check identity and gluability. \\
    The identity axiom is easy to verify as functions are determined by their values at points. Thus, if two functions agree on an open cover, they are the same function. \bbni
    For gluability, let $\{U_i\}$ be an open cover with locally constant functions $\{f_i\}$ that agree on overlaps. Let $f$ be the (set-theoretic) function that agrees with this information. We need to show that $f$ is locally constant. \\
    Let $p \in X$. Then there exists $i$ such that $p \in U_i$. Thus, there exists a neighborhood $W \subseteq U_i$ containing $p$ such that $f_i$ is constant on it. Then $f$ restricted to this is constant. Thus, $f$ is locally constant.
\end{solution}

\begin{exercise}{2.2.F.} 
    Let $Y$ be a topological space. To each open set $U$ of $X$, associate the set of continuous maps $U \to Y$. Show that this forms a sheaf. (2.2.D(b) and 2.2.E are special cases.) 
\end{exercise}
\begin{solution}
    We define the sheaf $\FF$ on $X$ as follows:
    \begin{enumerate}[itemsep=1pt, topsep=1pt]
        \item $\FF(U)$ is the set of continuous maps to $Y$
        \item For $V \subseteq U$, $res_{U,V}(f) = f|_V$, the usual restriction.
    \end{enumerate}
    This is clearly a presheaf via the usual properties of functions. We have to show the identity and gluability axioms. \\
    For the identity axiom, note that a continuous map is determined by the map on underlying sets. Thus, if two functions agree on an open cover, they agree on the underlying set, and thus constitute the same function. \\
    For the gluability axiom, let $\{U_i\}$ be an open cover of $X$ and $f_i \in \FF(U_i)$ be continuous functions to $Y$ such that $f_i|_{U_i \cap U_j} \equiv f_j|_{U_i \cap U_j}$. Let $f$ be the (unique) function of sets that agrees with this information. We need to show $f$ is continuous. \\
    Let $V \subseteq Y$ be open. Let $U = f^{-1}(V)$. Then, 
    \[ U = f^{-1}(V) = \bigcup_i f_i^{-1}(V) \]
    with each of the $f_i^{-1}(V)$ open (pre-image of an open set under a continuous function). Let $p \in U$ be arbitrary. Then, $p \in f_i^{-1}(V)$ for some $i$. Thus, $U$ is open. Thus, $f$ is continuous.
\end{solution}

\begin{exercise}{2.2.G.}
    \begin{enumerate}[itemsep=1pt, topsep=1pt]
        \item Suppose we are given a continuous map $\mu: Y \to X$. Show that ``section of $\mu$'' form a sheaf. That is, to each open set $U$ of $X$ associate the set of continuous maps $s: U \to Y$ such that $\mu \circ s = \id_U$. Show that this forms a sheaf.
        \item Suppose that $Y$ is a topological group. Show that continuous maps to $Y$ form a sheaf of groups.
    \end{enumerate}
\end{exercise}
\begin{solution}
    \todo{}
\end{solution}

\begin{exercise}{2.2.H.} 
    Suppose you have a map $\pi: X \to Y$ and a presheaf $\FF$ on $X$. Construct the pushforward presheaf on $Y$ by $\pi_*(V) := \OO(\pi^{-1}(V))$ for $V \subseteq Y$. Prove that this is a presheaf and is a sheaf if $\FF$ is.
\end{exercise}
\begin{solution}
    \todo{}
\end{solution}

\begin{exercise}{2.2.I.} 
    Show that the pushforward of sheaves induces a map on stalks $\pi_*(\FF)_{q} \to \FF_p$ where $q = \pi(p)$. 
\end{exercise}
\begin{solution}
    \todo{}
\end{solution}

\begin{exercise}{2.2.J.} 
    If $(X, \OO_X)$ is a ringed space and $\FF$ is a $\OO_X$-module, show that $\FF_p$ is a $\OO_{X, p}$-module.
\end{exercise}
\begin{solution}
    As $\FF$ is a $\OO_X$-module, we have the following commutative diagram for $V \subseteq U$:
    \[
    \xymatrix{
    \OO(U) \times \FF(U) \ar[r] \ar[d]_{res_{U,V} \times res_{U,V}} & \FF(U) \ar[d]^{res_{U,V}} \\
    \OO(V) \times \FF(V) \ar[r] &\FF(V)
    }
    \]
    Let $[(f, U)] \in \OO_{X,p}$ and $[(\chi, V)] \in \FF_p$. Define the action of $\OO_{X,p}$ by $[(f,U)] \cdot [(\chi, V)] = [(f \cdot \chi, U \cap V)]$. We need to show this is well-defined with respect to the equivalence relation on sections yielding the stalk. \\
    Suppose $[(f, U)] = [(g, U')] $. Then, there exists $W \subseteq U \cap U'$ such that $f|_W \equiv g|_W$. This also holds for $W \cap V$. Thus, $[(f \cdot \chi, U \cap V)] = [(g \cdot \chi, U' \cap V)]$ as the representatives agree over $W \cap V \subset U \cap U' \cap V$. An analogous argument holds for being well-defined with respect to the second entry.  
\end{solution}

\subsection{Morphisms:}
\begin{exercise}{2.3.A.} 
    If $\phi: \FF \to \GG$ is a morphism of sheaves, then $\phi$ induces a morphism of stalks $\phi_p: \FF_p \to \GG_p$. 
\end{exercise}
\begin{solution}
    Let $[(f, U)] \in \FF_p$. We define the induced map $\phi_p([(f, U)]) = [(\phi(U)(f), U)] \in \GG_p$. We need to show this is well-defined. \\
    If $(f, U) \simeq (g, V)$, then there exists open $W \subset U \cap V$, $p \in W$ such that $f|_W = g|_W$. By the definition of sheaf morphisms, this implies that:
    \[ \phi(U)(f)|_W = \phi(V)(g)|_W\]
    Thus, $[(\phi(U)(f), U)] = [(\phi(V)(g), V)]$. 
\end{solution}

\begin{exercise}{2.3.B.} 
    Suppose $\pi: X \to Y$ is continuous. Show that the pushforward gives a functor $\pi_*: Sets_X \to Sets_Y$. 
\end{exercise}
\begin{solution}
    Define the pushforward functor as follows: 
    \begin{enumerate}[topsep=1pt, itemsep=1pt]
        \item For $\FF \in Sets_X$, $\pi_*(\FF)(U) := \FF(\pi^{-1}(U))$ for $U \subseteq Y$ open.
        \item For $f: \FF \to \GG$ sheaf morphism with $\FF, \GG \in Sets_X$, define $\pi_*(f): \pi_*(\FF) \to \pi_*(\GG)$ by $\pi_*(f)(U) := f(\pi^{-1}(U))$ (this is a morphism from $\FF(\pi^{-1}(U))) \to \GG(\pi^{-1}(U))$.
    \end{enumerate}
    We need to show that $\pi_*$ preserves identities and compositions. \\
    Let $\FF$ be a sheaf on $X$ and $\id_\FF$ be the identity morphism. Then, we note that this preserves identities:
    \[ \pi_*(\id_\FF)(U) := \id_\FF(\pi^{-1}(U)) = \id_{\FF(\pi^{-1}(U))} = \id_{\pi_*(\FF)(U)}\] 
    Next, let $\FF, \GG, \HH$ be sheafs on $X$ with sheaf morphism $\phi: \FF \to \GG$ and $\psi: \GG \to \HH$. Let $f \in \pi_*(\FF)(U) = \FF(\pi^{-1}(U))$ be a section. Then, 
    \[\pi_*(\phi)(f) = \phi(\pi^{-1}(U))(f) \in \GG(\pi^{-1}(U)) = \pi_*(\GG)(U)\] 
    Similarly, 
    \[\pi_*(\psi)(\phi(\pi^{-1}(U))(f)) = \psi(\pi^{-1}(U))(\phi(\pi^{-1}(U))(f)) \in \HH(\pi^{-1}(U)) = \pi_*(\GG)(U)\]  
    Overall,  
    \begin{align*}
        \pi_*(\psi) \circ \pi_*(\phi)(f) &:= (\psi(\pi^{-1}(U)) \circ \phi(\pi^{-1}(U)))(f) = \pi_*(\psi \circ \phi)
    \end{align*}

\end{solution}

\begin{exercise}{2.3.C.} 
Suppose $\FF$ and $\GG$ are two sheaves of sets on $X$.
(In fact, it will suffice that $\FF$ is a presheaf.)
Let $\mathcal{H}om(\FF,\GG)$ be the collection of data
\[ \mathcal{H}om(\FF,\GG)(U) := \mathrm{Mor}(\FF|_U, \GG|_U) \]
Show that this is a sheaf of sets on $X$.
\end{exercise}
\begin{solution}
    We define the restriction maps naturally as a sheaf morphism contains data for a sheaf morphism on a subspace. This is clearly a presheaf of sets. To prove this is a sheaf, we verify the axioms:
    \begin{enumerate}[itemsep=1pt, topsep=1pt]
        \item (Identity). Let $\{U_i\}$ be an open cover of $U$. Let $\phi, \psi \in \HH om(\FF, \GG)(U)$ be two sections such that $\phi|_{U_i} = \psi|_{U_i}$ for all $i$. Let $W \subseteq U$ be some open set. It suffices to show that $\phi(W) = \psi(W)$ as maps. \\ 
        Note that $\{W \cap U_i\}$ cover $W$. Moreover, we have $\phi(W \cap U_i) = \psi(W \cap U_i)$ (as $\phi|_{U_i} = \psi|_{U_i}$). Let $f \in \FF(W)$. Then, 
        \[ \phi(W \cap U_i)(f|_{W \cap U_i}) = \psi(W \cap U_i)(f|_{W \cap U_i}) \in \GG(W \cap U_i)\]
        As $\GG$ is a sheaf, we get $\phi(W)(f) = \psi(W)(f)$ by the identity axiom. As $f$ was arbitrary, this implies $\phi(W) = \psi(W)$. We are done.  
        \item (Gluability). Let $\{U_i\}$ be an open cover of $U$. Let $\phi_i \in \HH om(\FF, \GG)(U_i)$ such that $\phi_i|_{U_i \cap U_j} = \phi_j|_{U_i \cap U_j}$. We need to construct $\phi \in \HH om(\FF, \GG)(U)$ which restricts to these. \\
        Let $W \subseteq U$ be open and $f \in \FF(W)$. Note that $\{ W \cap U_i\}$ covers $W$. Moreover,  
        \[ \phi_i(W \cap U_i)(f|_{W \cap U_i}) \in \GG(W \cap U_i)\]
        Next, note that: $\phi_i(W \cap U_i \cap U_j) = \phi_j(W \cap U_i \cap U_j)$. Thus, on $W \cap U_i \cap U_j$, we have:
        \[ \phi_i(W \cap U_i \cap U_j)(f|_{W \cap U_i \cap U_j}) = \phi_j(W \cap U_i \cap U_j)(f|_{W \cap U_i \cap U_j})\]
        By gluability on $\GG$, we get a section in $\GG(U)$, call it $g$, such that $g|_{W \cap U_i} = \phi_i(W \cap U_i)(f|_{W \cap U_i})$. Let $\phi(W)(f) = g$. This defines $\phi \in \HH om(\FF, \GG)(U)$ as we have defined $\phi$ for every open $W \subseteq U$. Moreover, we note that this restricts appropriately as for $W \subseteq U_i$ our definition collapses to $\phi(W \cap U_i) = \phi_i(W \cap U_i) = \phi_i(W)$. 
    \end{enumerate} 
\end{solution}

\newpage 
\subsection{Abelian Category of Presheaves:}

\begin{exercise}{2.3.E.} 
Show that $\ker_{{pre}} \Phi$ is a presheaf. (Hint: if $V \hookrightarrow U$, define the restriction map by chasing the following diagram:
\[ \xymatrix{
    0 \ar[r] & \ker_{pre} \Phi(U) \ar[r] \ar@{-->}[d]^{\exists !} & \FF(U) \ar[r] \ar[d]^{res_{U,V}} & \GG(U) \ar[d]^{{res}_{U,V}} \\
    0 \ar[r] & \ker_{{pre}} \Phi(V) \ar[r] & \FF(V) \ar[r] & \GG(V) } \]
You should check that the restriction maps compose as desired.
\end{exercise}
\begin{solution}
    Define the restriction map using the universal property of $\ker(\Phi(V))$. The uniqueness of this morphism implies that the restriction maps are compatible. 
\end{solution}

\begin{exercise}{2.3.F.} 

\end{exercise}
\begin{solution}
    \todo{}
\end{solution}


\begin{exercise}{2.3.I.} 

\end{exercise}
\begin{solution}
    \todo{}
\end{solution}

\begin{exercise}{2.3.J.} 
Let $X = \mathbb{C}$ with the classical topology, let $\underline{\mathbb{Z}}$ be the constant sheaf on $X$ associated to $\mathbb{Z}$, let $\mathcal{O}_X$ be the sheaf of holomorphic functions, and let $\FF$ be the \emph{presheaf} of functions admitting a holomorphic logarithm. Describe an exact sequence of presheaves on $X$:
\[ 0 \longrightarrow \underline{\mathbb{Z}} \longrightarrow \mathcal{O}_X \longrightarrow \FF \longrightarrow 0 \]
where $\underline{\mathbb{Z}} \to \mathcal{O}_X$ is the natural inclusion and $\mathcal{O}_X \to \FF$ is given by
\[ f \longmapsto \exp(2\pi i f) \]
(Be sure to verify exactness.) Show that $\FF$ is \emph{not} a sheaf.
(\emph{Hint:} $\FF$ does not satisfy the gluing axiom. The problem is that there are functions that do not have a logarithm but locally have a logarithm.)
\end{exercise}
\begin{solution}
    \todo{}
\end{solution}
\newpage
\subsection{Stalks:}

\begin{exercise}{2.4.A.} 
Prove that a section of a sheaf of sets is determined by its germs, i.e., the natural map
\[\FF(U) \longrightarrow \prod_{p \in U} \FF_p \]
is injective.
\end{exercise}
\begin{solution}
    Let $f, g \in \FF(U)$ be two sections whose germs agree at every point, i.e. for all $p \in U$, $[(f, U)] = [(g, U)]$. Then, there exists a $W_p \subseteq U$ such that $f|_{W_p} = g|_{W_p}$. Note that $\{W_p\}_{p \in U}$ form an open cover of $U$. By the identity axiom, we are done.
\end{solution}

\begin{exercise}{2.4.B.} 
    Show that the support of section $\mathrm{Supp}(s)$ is closed in $X$. 
\end{exercise}
\begin{solution}
    We will show that the complement of the support is open. Let $p \in X \setminus \mathrm{Supp}(s)$. Then, $s_p = 0 \in \FF_p$. This implies that there exists an open set $W_p$ containing $p$ such that $s|_{W_p} = 0|_{W_p}$, i.e. $W_p \subseteq X \setminus \mathrm{Supp}(s)$. Thus, $X \setminus \mathrm{Supp}(s)$ is open. 
\end{solution}

\begin{exercise}{2.4.C.} 
    Prove that any choice of compatible germs for a sheaf of sets $\FF$ over $U$ is the image of a section of $\FF$ over U. (Hint: you will use gluability.)
\end{exercise}
\begin{solution}
    Let $\{U_i\}$ be a cover of $U$ and $f_i \in \FF(U_i)$ be functions such that the germ of $f_i$ at $p \in U_i$ is $s_p$. That is, $s_p$ for $p \in U$ are compatible germs. First note that these sections agree on overlaps as they agree at each stalk in the intersection (2.4.A). Thus, by gluability, we are done. 
    
\end{solution}

\begin{exercise}{2.4.D.} 
    If $\phi_1$ and $\phi_2$ are morphisms from a presheaf of sets $\FF$
    to a sheaf of sets $\GG$ that induce the same maps on each stalk,
    show that $\phi_1 = \phi_2$. Hint: consider the following diagram.
    \[ \xymatrix{
    \FF(U) \ar[r] \ar[d] &
    \GG(U) \ar[d] \\
    \displaystyle\prod_{p \in U} \FF_p \ar[r] &
    \displaystyle\prod_{p \in U} \GG_p} \]
\end{exercise}
\begin{solution}
    Assume $\phi_1$ and $\phi_2$ induce the same maps on stalks. Let $U \in X$ be an arbitrary open. It suffices to show that $\phi_1(U) = \phi_2(U)$. Let $f \in \FF(U)$. Consider the map $\FF(U) \to \prod_{p\in U} \FF_p \to \prod_{p \in U} \GG_p$. Since $\phi_1$ and $\phi_2$ induce the same maps on stalks, this map is identical regardless of whether the top morphism is $\phi_1$ or $\phi_2$. Thus, the germs of $\phi_1(U)(f)$ and $\phi_2(U)(f)$ agree for all $p \in U$. Then, as $\GG$ is a sheaf, we use $(2.4.A)$ to conclude $\phi_1(U)(f) = \phi_2(U)(f)$. As $f$ was arbitrary, we are done. 
\end{solution}

\begin{exercise}{2.4.E.} 
Show that a morphism of sheaves of sets is an isomorphism if and only if
it induces an isomorphism of all stalks.
Hint: Use $2.4.4.1$. Once you have injectivity, show
surjectivity, perhaps using Exercise 2.4.C, or gluability in some other
way; this is more subtle.
\end{exercise}
\begin{solution}
     \todo{}
\end{solution}

\begin{exercise}{2.4.F.} 

\end{exercise}
\begin{solution}
    \todo{}

\end{solution}

\subsection{Sheafification:}
\begin{exercise}{2.4.H.} 
    Show that for any morphism of pre-sheaves $\phi: \FF \to \GG$ we have an induced map of sheaves $\phi^{sh}: \FF^{sh} \to \GG^{sh}$. Show that sheafification is a functor. 
\end{exercise}
\begin{solution}
    Consider the following diagram: 
    \[ \xymatrix{ \FF \ar[r]^{\phi} \ar[d]_{sh} & \GG \ar[d]^{sh} \\
    \FF^{sh} \ar@{-->}[r]_{\phi^{sh}} & \GG^{sh} }\]
    where the map $\phi^{sh}$ is induced by the universal property of $\GG^{sh}$. Functoriality is easy to see given this commutative diagram.
\end{solution}

\begin{exercise}{2.4.I.} 
    Show that $\FF^{sh}$ forms a sheaf. 
\end{exercise}
\begin{solution}
    Take the restriction map to be $(f_p)_{p \in U} \to (f_p)_{p \in V}$ for $V \subseteq U$. Note that the image also consists of compatible germs by taking appropriate intersections with $V$ in the definition of a section. \\
    To show identity, we let $\{U_i\}$ be an open cover of $U$ with $f, g \in \FF^{sh}(U)$ such that $f|_{U_i} = g|_{U_i}$ for all $i$. Then, $(f_p)_{p \in U_i} = (g_p)_{p \in U_i}$. Thus, $(f_p)_{p \in U} = (g_p)_{p \in U}$ as $U_i$ cover. Thus, $f = g$. \\ 
    To show gluability, let $f_i \in \{U_i\}$ be sections such that $f_i|_{U_{i} \cap U_{j}} = f_j|_{U_{i} \cap U_{j}}$. Note that for $p \in U$, we can pick a germ at $p$ by taking $U_i \ni p$ and taking $(f_i)_p$. This is well-defined by the agreement on overlaps, i.e. $(f_i)_{p} = (f_j)_p$ for $p \in U_i \cap U_j$. These germs are compatible by construction and thus define a section in $\FF(U)$.
\end{solution}

\begin{exercise}{2.4.J.} 
    Describe a natural map of presheaves $\FF \to \FF^{sh}$. 
\end{exercise}
\begin{solution}
    Let $\phi(U): \FF(U) \to \FF^{sh}(U)$ be $f \to (f_p)_{p \in U}$. This is a map of presheaves as it commutes with restrictions. That is, for $V \subseteq U$, we have $(f|_V)_p = f_p$ as they agree on $V \subseteq U$. 
\end{solution}

\begin{exercise}{2.4.K.} 
    Show that $\FF^{sh}$ satisfies the universal property of sheafification. 
\end{exercise}
\begin{solution}
    Let $\phi: \FF \to \GG$ be a map of presheaves with $\GG$ a sheaf. We first construct the map $\psi: \FF^{sh} \to \GG$. \\
    By $(2.3.A)$, $\phi$ induces a morphism on stalks, $\phi_p: \FF_p \to \GG_p$ with $\phi_p(f_p) = \phi(f)_p$. First, we claim that if $(f_p)_{p \in U}$ with $f_p \in \FF_p$ are compatible germs, then $(\phi_p(f_p))_{p \in U}$ are also compatible. Clearly, if for all p there exists $V \subseteq U$ containing $p$ and $s \in \FF(V)$ such that $s_q = f_q$ for all $q \in V$. But then, $\phi(s)_q = \phi_q(s_q) = \phi_q(f_q)$ with $\phi(s) \in \GG(V)$. Next, using $(2.4.A)$ and $(2.4.C)$, we note that sections of a sheaf are uniquely determined by compatible germs. Thus, we have $g \in \GG(U)$ with $\phi(f_p) = g_p$. Thus, we get a maps $\psi(U): \FF^{sh}(U) \to \GG(U)$ by $(f_p)_{p \in U} \mapsto g$ with $g_p = \phi_p(f_p)$. \bbni    
    Clearly, $\psi$ defines a map of sheaves as restriction commutes with $\psi$, i.e. for $V \subseteq U$, $g|_V$ satisfies $(g|_V)_p = \phi_p(f_p)$ for $p \in V$ and by $(2.4.A)$ is unique. Thus, we only need to show $\phi = \psi \circ (\cdot)^{sh}$. \bbni 
    Let $f \in \FF(U)$. Then, 
    \[ \psi(U) \circ (\cdot)^{sh}(U)(f) = \psi(U)((f_p)_{p \in U}) = g \]
    with $g_p = \phi_p(f_p)$ for all $p \in U$. We note that $\phi(U)(f)$ satisfies this property, so my uniqueness, $\phi(U)(f) = g$. Hence, we win. 
\end{solution}

\begin{exercise}{2.4.L.}

\end{exercise}
\begin{solution}
    \todo{}
\end{solution}

\begin{exercise}{2.4.M.} 

\end{exercise}
\begin{solution}
    \todo{}

\end{solution}

\subsection{Subsheaves and Quotient Sheaves:}
\begin{exercise}{2.4.N.} 

\end{exercise}
\begin{solution}
    \todo{}

\end{solution}

\begin{exercise}{2.4.O.} 

\end{exercise}
\begin{solution}
    \todo{}

\end{solution}

\begin{exercise}{2.4.P.} 

\end{exercise}
\begin{solution}
    \todo{}

\end{solution}
\subsection{Sheaf on a base:}


\begin{exercise}{2.5.A.} 
    Let $\{B_i\}$ is the base of a topology on $X$. You are given the information $\{\FF(B_i), res_{B_i, B_j} \}$. Show that you can recover the entire sheaf from this.
\end{exercise}
\begin{solution}
    \todo{}
    % We will show that you can recover germs and their compatibility through this information. \\
    % Take $[(f, W)] \in \FF_p$. Then there exists a $B_i \subseteq W$ containing $p$. Then $(res_{W, B_i}(f), B_i)$ represents the same germ as $(f, W)$. Thus, we can recover all germs using just the sheaf on the base. \bbni 
    % Similarly, if $[(f, U)] = [(g, V)]$ are compatible germs then there exists $W \subseteq U \cap V$ with $f|_W \equiv g|_W$. Then, there exists $B_i \subseteq W$ containing $p$ where this equivalence also holds. Thus, the sheaf on a base determines the equivalence of germs. \\
    % Thus, as this information is sufficient to determine the stalks, it is sufficient to determine the sheaf.
\end{solution}

\begin{exercise}{2.5.B.}    
    Verify that $F(B) \to \FF(B)$ is an isomorphism, likely by showing that it is injective and surjective (or else by describing the inverse map and verifying that it is indeed inverse). (Hint: elements of $\FF(B)$ are determined by stalks, as are elements of $F(B)$.)
\end{exercise}
\begin{solution}
    \todo{}
\end{solution}

\begin{exercise}{2.5.C.}
    { \textbf{(Morphisms of Sheaves $\leftrightarrow$ morphisms of sheaves on a base)}} \\
    Suppose $\{B_i\}$ is a base for the topology on $X$. Suppose you are given a sheaf morphism $\FF \to \GG$ on the base, i.e., you have for $B_j \subseteq B_i$: 
    \[ 
        \xymatrix{ \FF(B_i) \ar[r] \ar[d]_{res_{B_i, B_j}} &\GG(B_i) \ar[d]^{res_{B_i, B_j}} \\ 
        \FF(B_j) \ar[r] &\GG(B_j)}
    \]
    Show that: 
    \begin{enumerate}[topsep=1pt, itemsep=1pt]
        \item Verify that a morphism of sheaves is determined by the induced morphism of sheaves on a base. 
        \item Show that morphism of sheaves on a base gives a morphism of induced sheaves. 
    \end{enumerate}
\end{exercise}
\begin{solution}
    \todo{}

\end{solution}

\begin{exercise}{2.5.D.} 
    Suppose $X = \bigcup U_i$ is an open cover of $X$ and we have sheaves $\FF_i$ on $U_i$ along with isomorphism $\phi_{ij}: \FF_i|_{U_i \cap U_j} \to \FF_j|_{U_i \cap U_j}$ that agree on triple overlaps, i.e. $\phi_{jk} \circ \phi_{ij} = \phi_{ik}$ (cocycle condition). Show that these can be glued together into a sheaf $\FF$ on $X$ which is unique up to unique isomorphism with obvious restrictions. 
\end{exercise}
\begin{solution}
    \todo{}
\end{solution}
